\chapter{Espacios vectoriales normados}
Durante este capítulo se asume que $V$ es un espacio vectorial sobre $\R$ .
\section{Normas}
Se abstrae la noción de distancia a través la norma.
\begin{definition}\label{def:norma}
  Llamamos norma en $V$ a una función $\norm{\cdot} : V\to \R$ si cumple que para todo $v,w\in V$ y $\lambda \in \R$
  \begin{enumerate}
    \item $\norm{x}\ge 0$ y $\norm{x}=0$ si y solo si $x=0$. (\textit{definida positiva})
    \item $| \lambda|\norm{v}=\norm{\lambda v}$  (\textit{homogeneidad}).
    \item $\norm{v+w}\le \norm{v}+\norm{w}$ (\textit{desigualdad triangular}).
  \end{enumerate}
\end{definition}
\begin{lemma}[Desigualdad triangular inversa]
  Sean $v,w\in V$, entonces $|\norm{v}-\norm{w}|\le \norm{v-w}$
\end{lemma}
\begin{proof}
  Véase que $\norm{v}=\norm{(v-w)+w}$, aplicamos la desigualdad triangular $\norm{v}\le \norm{v-w}+\norm{w}$ y $\norm{v}-\norm{w}\le \norm{v-w}$; de forma análoga $\norm{w}-\norm{v}\le \norm{v-w}$.
\end{proof}
Las siguientes normas son las más utilizadas en $\R^n$.
\begin{definition}[Normas $L_p$]
  Sea $p\ge 1$ y $V=R^n$, llamamos norma $L_p$ a  
  \[
    \norm{x}_p=\sqrt[p]{\abs{x_1}^p+\ldots+\abs{x_n}^p}
  \] donde $x=(x_1,\ldots,x_n)$.\\
  Además definimos la p-infinita por $\norm{x}_\infty=\max_{1\le i \le n} \{\abs{x_i}\}$ 
\end{definition}
Si son normas han de cumplir la desigualdad triangular; para demostrar esto usaremos una serie de lemas, cada uno demostrado por el anterior.
\begin{definition}[Función convexa]
  Llamamos a la función $f:[a,b]\subset \R\to \R$ convexa si $f''(x)\ge 0$.
\end{definition}
Entonces para $x<\alpha<y$ tenemos que \[
    \frac{f(\alpha))-f(x))}{\alpha-x}\le \frac{f(z))-f(\alpha))}{z-\alpha}.
  \] Ademas sabemos que para $t\in[0,1]$, $x<tx+(1-t)y<z$, entonces sustituyendo $\alpha=tx+(1-t)y$ y haciendo las cuentas (creeme) llegamos a que \[
      f(tx+(1-t)y)\le tf(x)+(1-t)f(y)
  \] 
\begin{lemma}\label{lemma:convexa-contractiva}
  Sean $\alpha ',\beta '\geq 0$ tal que $\alpha'+\beta'=1$ , si $f:D\subseteq \R \to \R$ es una función convexa, entonces para todo $x,y\in D$ \[
      \alpha' f(x)+\beta' f(y)\ge f(\alpha' x+\beta' y)
  \]
\end{lemma}
\begin{lemma}[Desigualdad de Young]\label{lemma:young}
  Sean $a,b\ge 0$ y $\alpha,\beta>1$ tal que $\frac{1}{\alpha}+\frac{1}{\beta}=1$, entonces 
  \[
      ab\le \frac{a^\alpha}{\alpha}+\frac{b^\beta}{\beta}
  \]  
\end{lemma}
Atentos a la magia negra.
\begin{proof}
  Véase que $f(x)=\exp(x)$ es convexa en su dominio. Usando el lema \ref{lemma:convexa-contractiva} tenemos que \[
      \frac{1}{\alpha}\exp(x)+\frac{1}{\beta}\exp(y)\ge \exp(\frac{1}{\alpha}x+\frac{1}{\beta}y)
  \] 
  eligiendo $x=\log(a^\alpha)$ e $y=\log(b^\beta)$ y simplificando, entonces 
  \[
      \frac{a^\alpha}{\alpha}+\frac{b^\beta}{\beta}\ge ab
  \]  
\end{proof}
\begin{lemma}[Desigualdad de Hölder]
  Sean $x,y\in \R^n$, tal que $x_{i},y_{i}\ge 0$ para $i\in \{1,\ldots,n\}$; $\alpha,\beta>1 $ tal que $\frac{1}{\alpha}+\frac{1}{\beta}=1$, entonces 
  \[
    \sum_{i=1}^{n} x_{i}y_{i}\le \norm{x}_\alpha\cdot\norm{y}_\beta
  \] 
\end{lemma}
\begin{proof}
  Si $\norm{x}_\alpha=0$ o $\norm{y}_\beta=0$, entonces es trivial; asume que no lo son. Basta demostrar que $\frac{\sum_{i=1}^{n} x_{i} y_{i}}{\norm{x}_\alpha \cdot\norm{y}_\beta}\le 1$. Para esto, consideramos el $i$-ésimo termino de la expresión anterior y aplicamos la desigualdad de Young \ref{lemma:young} \[
    \parens*{ \frac{x_{i}}{\norm{x}_\alpha}}\cdot 
    \parens*{\frac{y_{i}}{\norm{y}_\beta}}\le 
\frac{1}{\alpha}\parens*{ \frac{x_{i}^\alpha}{\norm{x}_\alpha^\alpha}}+
    \frac{1}{\beta}\parens*{\frac{y_{i}^\beta}{\norm{y}_\beta^\beta}} 
  ;\]
  sumamos ahora los $n$ términos y agrupamos \[
     \sum_{i=1}^{n} \frac{x_{i}y_{i}}{\norm{x}_\alpha\norm{y}_\beta}\le 
\frac{1}{\alpha}\parens*{ \frac{\sum_{i=1}^{n} x_{i}^\alpha}{\norm{x}_\alpha^\alpha}}+
    \frac{1}{\beta}\parens*{\frac{\sum_{i=1}^{n} y_{i}^\beta}{\norm{y}_\beta^\beta}}=\frac{1}{\alpha}+\frac{1}{\beta}=1
  \] 
\end{proof}
\begin{lemma}[Desigualdad de Minkowski]
  Sean $x,y\in\R^n$ y $p\ge 1$, entonces \[
      \norm{x+y}_p\le \norm{x}_p+\norm{y}_p
  \] 
\end{lemma}
\begin{proof}
  Si $p=1$, sumando $|x_k+y_k|\le |x_k|+|y_k|$ obtenemos $\norm{x+y}_1\le \norm{x}_1+\norm{y}_1$.\\
  Sea $p>1$ y $q=\frac{p}{p-1}$, de modo que $\frac{1}{p}+\frac{1}{q}=1$. Denotamos $a_k=|x_k|$ y $b_k=|y_k|$; como $|x_k+y_k|\le a_k+b_k$ se tiene $\norm{x+y}_p\le \parens*{\sum_{k=1}^n(a_k+b_k)^p}^{1/p}$ y $\norm{a}_p=\norm{x}_p$, $\norm{b}_p=\norm{y}_p$, luego basta demostrar que \[
    \parens*{\sum_{k=1}^n(a_k+b_k)^p}^{1/p}\le \norm{a}_p+\norm{b}_p.
  \]
  Escribimos
  \[
    \sum_{k=1}^n (a_k+b_k)^p = \sum_{k=1}^n a_k(a_k+b_k)^{p-1}+\sum_{k=1}^n b_k(a_k+b_k)^{p-1}
  \]
  y aplicamos Hölder a cada suma con exponentes $p$ y $q$. Como $(p-1)q=(p-1)\cdot\frac{p}{p-1}=p$,
  \begin{align*}
    \sum_{k=1}^n a_k(a_k+b_k)^{p-1}&\le \norm{a}_p\parens*{\sum_{k=1}^n(a_k+b_k)^p}^{1/q},\\
    \sum_{k=1}^n b_k(a_k+b_k)^{p-1}&\le \norm{b}_p\parens*{\sum_{k=1}^n(a_k+b_k)^p}^{1/q};
  \end{align*}
  sumando,
  \[
    \sum_{k=1}^n(a_k+b_k)^p\le \parens*{\norm{a}_p+\norm{b}_p}\parens*{\sum_{k=1}^n(a_k+b_k)^p}^{1/q}.
  \]
  Si $\sum_{k=1}^n(a_k+b_k)^p=0$ la desigualdad es trivialmente $0\le 0$. En caso contrario dividimos ambos lados por $\parens*{\sum_{k=1}^n(a_k+b_k)^p}^{1/q}$; como $1-\frac{1}{q}=\frac{1}{p}$ obtenemos
  \[
    \parens*{\sum_{k=1}^n(a_k+b_k)^p}^{1/p}\le \norm{a}_p+\norm{b}_p.
  \]
\end{proof}
\begin{theorem}[Las normas $L_p$ son normas]
\end{theorem}
\begin{definition}[Normas equivalentes]
  Llamamos a dos normas $\norm{\cdot}, \norm{\cdot}'$ equivlaentes si existen $C_1,C_2>0$ tal que para todo $x\in V$ \[
      C_1 \norm{x}\le \norm{x}' \le C_2\norm{x}
  \] 
\end{definition}
\begin{theorem}[Equivalencia de normas en $\R^n$ ]\label{thm:equivalencia_de_normas_en_Rn}
  En $\R^n$ todas las normas son equivalentes
\end{theorem}
\begin{proof}
  De momento esto te lo vas a tener que creer.
\end{proof}
\section{Producto interno}
\begin{definition}[Producto interno]
  Llmamamos producto interno en $V$ a una función bilineal  simétrica definida positiva $\ip{x,y}: V\times V\to \R$ \begin{enumerate}
    \item $\ip{x,x}\ge 0$ dandose la igualdad si y solo si $x=0$,
    \item para todo $y\in V$ la función $x\mapsto \ip{x,y}$ es lineal,
    \item $\ip{x,y} = \ip{y,x}$.
  \end{enumerate}
\end{definition}
\begin{proposition}[Norma inducida por un producto interno]
  Dado un producto interno, la función $\norm{x}=\sqrt{\ip{x,x}} $ es una norma.
\end{proposition}
Para demostrar que cumple la desigualdad triangular usaremos el \fullref{thm:cauchy-schwarz}, y para demostrar este el siguiente lema.
\begin{lemma}\label{lemma:discriminante}
  Sea $p(t)=at^2+bt+c$, con $a,b,c,t\in \R$, tal que $p(t)\geq 0$ para todo $t\in \R$, entonces el discriminante es no positivo \[
      \Delta=b^2-4ac \leq 0
  \] 
\end{lemma}
La idea es sencilla: como $a>0$, la parábola $p(t)$ abre hacia arriba. Si $p(t)\ge 0$ para todo $t$, entonces nunca cruza el eje horizontal, luego tiene a lo sumo una raíz (la del vértice) y por tanto $\Delta\le 0$.
\begin{center}
\begin{tikzpicture}[scale=0.75, line cap=round, >=Latex]
  % Left: Delta > 0, parabola crosses axis
  \begin{scope}[shift={(-4,0)}]
    \draw[->] (-2.2,0) -- (2.5,0);
    \draw[->] (0,-0.8) -- (0,2.8);
    \draw[thick, blue, domain=-1.8:2.1, samples=60] plot (\x, {0.7*(\x+0.8)*(\x-1.5)});
    \fill[blue] (-0.8,0) circle (2pt);
    \fill[blue] (1.5,0) circle (2pt);
    \node[below] at (0.35,-0.9) {$\Delta>0$};
  \end{scope}
  % Right: Delta <= 0, parabola stays non-negative
  \begin{scope}[shift={(4,0)}]
    \draw[->] (-2.2,0) -- (2.5,0);
    \draw[->] (0,-0.8) -- (0,2.8);
    \draw[thick, red, domain=-1.8:2.1, samples=60] plot (\x, {0.6*(\x-0.2)*(\x-0.2)+0.4});
    \node[below] at (0.35,-0.9) {$\Delta\le 0$};
    \node[right, red] at (1.8,2.2) {$p(t)\ge 0$};
  \end{scope}
\end{tikzpicture}
\end{center}

\begin{proof}
  Como $p(t)\ge 0$ para todo $t$ y $\lim_{t\to \infty}p(t)=\operatorname{signo}(a)\cdot\infty$, necesariamente $a>0$. Evaluamos $p$ en el vértice $t_0=-\frac{b}{2a}$: \[
    0\le p(t_0)=a\frac{b^2}{4a^2}-\frac{b^2}{2a}+c=\frac{4ac-b^2}{4a}=\frac{-\Delta}{4a}.
  \]
  Como $a>0$, multiplicando por $4a$ obtenemos $-\Delta\ge 0$.
\end{proof}
\begin{theorem}[Cauchy-Schwarz]\label{thm:cauchy-schwarz}
  Sea $x,y\in V$, si la norma es la inducida por el producto interno, entonces \[
    \abs{\ip{x,y}}\le \norm{x}\norm{y}
  \] 
\end{theorem}
\begin{proof}
  Como la norma está inducida por el producto interno, para todo $t\in \R$ tenemos \[
    0\le \norm{x+ty}^2=\ip{x+ty,x+ty}=\norm{y}^2t^2+2\ip{x,y}t+\norm{x}^2.
  \]
  Este es un polinomio de grado 2 en $t$ que es no negativo para todo $t$; por el \cref{lemma:discriminante} su discriminante es no positivo: \[
    \Delta=4\ip{x,y}^2-4\norm{x}^2\norm{y}^2\le 0,
  \]
  dividiendo por $4$ obtenemos $\ip{x,y}^2\le \norm{x}^2\norm{y}^2$, y tomando raíces concluimos $\abs{\ip{x,y}}\le \norm{x}\norm{y}$.
\end{proof}
Verifica por ti mismo que ahora se puede demostrar que es una norma.
\begin{proposition}[Ley del paralelogramo]
  Una norma está inducida por un producto interno si y solo si cumple que para todo $x,y\in V$ \[
    \norm{x+y}^2+\norm{x-y}^2=2\norm{x}^2+2\norm{y}^2.
  \] 
\end{proposition}
\begin{proof}
  Asume que la norma $\norm{\cdot}$ está inducida por un producto escalar, observa que $\norm{x\pm y}^2=\ip{x\pm y,x\pm y}=\norm{x}^2+\norm{y}^2\pm 2\ip{x,y}$ y entonces se verifica la ley del paralelogramo.\\
  Se deja como tortura para el lector verificar que el enunciado converso es cierto. Es dificil, no lo intentes.
\end{proof}
En $\R^2$ se ilustra de la siguiente forma y se puede demostrar usando geometría elemental que la ley se da en todo paralelogramo. 
% Requires: \usepackage{tikz}
% Optional but nice: \usetikzlibrary{arrows.meta,calc}

\begin{figure}[ht]
\centering
\begin{tikzpicture}[scale=1.05, line cap=round, line join=round, >=Latex]
  % Vectors
  \coordinate (O) at (0,0);
  \coordinate (x) at (3,1);
  \coordinate (y) at (1.2,2.6);
  \coordinate (xpy) at ($(x)+(y)$);

  % Axes (optional)
  \draw[->, gray!60] (-0.4,0) -- (4.8,0) node[below] {$e_1$};
  \draw[->, gray!60] (0,-0.4) -- (0,4.6) node[left] {$e_2$};

  % Parallelogram edges
  \draw[thick] (O) -- (x) -- (xpy) -- (y) -- cycle;

  % Vectors x and y
  \draw[very thick, ->] (O) -- (x) node[midway, below right] {$x$};
  \draw[very thick, ->] (O) -- (y) node[midway, left] {$y$};

  % Diagonals: x+y and x-y
  \draw[very thick, ->, blue] (O) -- (xpy) node[midway, above right] {$x+y$};
  \draw[very thick, ->, red]  (y) -- (x) node[midway, above left] {$x-y$};

  % Point labels
  \fill (O) circle (1.2pt) node[below left] {$0$};
  \fill (x) circle (1.2pt);
  \fill (y) circle (1.2pt);
  \fill (xpy) circle (1.2pt);

\end{tikzpicture}
\end{figure}

\section{Bolas, segmentos y conjuntos convexos}
Habiendo generalizado la noción de distancia, generalizamos la de intervalo.
\begin{definition}[Bola abierta]
  Definimos la bola abierta con centro $a\in V$ y radio $\epsilon>0$ \[
      B_\epsilon (a)=\{v\in V : \norm{v-a}<\epsilon\} 
  \]
\end{definition}
\begin{example}
  Sea $V=\R^2$, ilustramos las bolas abiertas con centro 0 y radio 1 para las normas $L_p$, $p=1,2,3,4,\infty$ en la \cref{fig:tema_1_bolas_lp}.
\end{example}
% fig/lp_balls.tex
\begin{figure}[ht]
\centering
\begin{tikzpicture}

% local helpers (figure-specific)
\newcommand{\LpBoundary}[1]{%
  \addplot[domain=0:360,samples=361,very thick]
  ({cos(x)/pow(pow(abs(cos(x)),#1)+pow(abs(sin(x)),#1),1/#1)},
   {sin(x)/pow(pow(abs(cos(x)),#1)+pow(abs(sin(x)),#1),1/#1)});
}
\newcommand{\LpFill}[1]{%
  \addplot[domain=0:360,samples=361,draw=none,fill opacity=0.12]
  ({cos(x)/pow(pow(abs(cos(x)),#1)+pow(abs(sin(x)),#1),1/#1)},
   {sin(x)/pow(pow(abs(cos(x)),#1)+pow(abs(sin(x)),#1),1/#1)})
  \closedcycle;
}
\newcommand{\SamplePoints}{%
  \addplot[only marks, mark=*, mark size=1.6pt] coordinates {(0.3,0.4)} node[above right] {$m_1$};
  \addplot[only marks, mark=*, mark size=1.6pt] coordinates {(1,0)} node[below right] {$m_2$};
  \addplot[only marks, mark=*, mark size=1.6pt] coordinates {(0,1)} node[above left] {$m_3$};
}

\begin{groupplot}[
  group style={group size=3 by 2, horizontal sep=1.2cm, vertical sep=1.2cm},
  width=5.2cm, height=5.2cm,
]

\nextgroupplot[title={$p=1$}, myBallAxis]
  \LpFill{1}\LpBoundary{1}\SamplePoints

\nextgroupplot[title={$p=2$}, myBallAxis]
  \LpFill{2}\LpBoundary{2}\SamplePoints

\nextgroupplot[title={$p=3$}, myBallAxis]
  \LpFill{3}\LpBoundary{3}\SamplePoints

\nextgroupplot[title={$p=4$}, myBallAxis]
  \LpFill{4}\LpBoundary{4}\SamplePoints

\nextgroupplot[title={$p=\infty$}, myBallAxis]
  \addplot[draw=none, fill opacity=0.12] coordinates {(-1,-1) (1,-1) (1,1) (-1,1)} \closedcycle;
  \addplot[very thick] coordinates {(-1,-1) (1,-1) (1,1) (-1,1) (-1,-1)};
  \SamplePoints

\nextgroupplot[title={}]
  \node[align=left] at (axis cs:0,0.4)
  {\(\displaystyle B_p(0,1)=\{(x,y):|x|^p+|y|^p<1\}\)\\[2mm]
   \(\displaystyle \overline B_p(0,1)=\{(x,y):|x|^p+|y|^p\le 1\}\)\\[2mm]
   \(\displaystyle B_\infty(0,1)=\{(x,y):\max(|x|,|y|)<1\}\).};

\end{groupplot}
\end{tikzpicture}
\caption{Bolas abiertas de radio \(1\) para \(\|\cdot\|_p\) en \(\mathbb R^2\), con puntos \(m_i\in \overline B_p(0,1)\).}
\label{fig:tema_1_bolas_lp}
\end{figure}

Generalizamos ahora la noción de segmento entre dos puntos.
\begin{definition}[Segmento]
  Dados dos puntos $x,y\in V$ llamamos segmento a \[
    [x,y]=\{z\in V : z=tx+(1-t)y, \;  t\in[0,1]\} 
  \] 
\end{definition}
\begin{definition}[Conjunto convexo]
  Decimos que un conjunto $C\subseteq V$ es convexo si para todo $x,y\in C$ el segmento que los une está contenido $[x,y]\subseteq  C$ 
\end{definition}
\begin{lemma}[Toda bola abierta es un conjunto convexo]
\end{lemma}
 \begin{proof}
   Sea $x,y\in B_\epsilon (a)$ y $t\in[0,1]$, entonces
   \begin{align*}
     \norm{[tx+(1-t)y]-a}=\norm{[tx+(1-t)y]-a+ta-ta}\\
     =\norm{t(x-a)+(1-t)(y-a)}\leq t\norm{x-a}+(1-t)\norm{y-a}<\epsilon
   \end{align*}
   y entonces $tx+(1-t)y\in B_\epsilon(a)$ 
 \end{proof}
\begin{definition}[Función convexa generlizada]
  Sea $C\subseteq V$ convexo, llamamos a $f: C\to \R$ convexa si para todo $x,y\in C$ y $t\in [0,1]$ \[
    f(tx+(1-t)y)\leq tf(x)+ (1-t)f(y)
  \] 
\end{definition}
Observa que toda norma es una función convexa, por la desigualdad triangular.
\section{Suceesiones}
\begin{definition}
  Denotamos con $(x^n)$ a la sucesión de terminos donde $x^k\in V$ para $k\in \N$ 
\end{definition}
Vease que usamos la notación $x^k$ para el $k$-ésimo elemento de la sucesión, mientras que $x_k$ denota la entrada.
\begin{definition}[Convergencia]
  La sucesión $(x_n)$ converge a $L\in V$ si \[
    \lim_{n\to\infty}\norm{x-L}=0
  \], donde este es un límite en $\R$.
\end{definition}
Véase a que esto equivale que para todo $\epsilon>0$ existe $n_0\in \N$ tal que si $k\geq n_0$, entonces $x_k\in B_\epsilon (L)$
\begin{lemma}[Propiedades de las suceesiones]
Si $(x_n)\to x$ e $(y_n)\to y$, entonces 
\begin{itemize}
  \item $(\norm{x_k})\to \norm{x}$ 
  \item para una sucesión real $(\lambda^k)\to \lambda$, $(\lambda^k x^k+y^k)\to \lambda x+y$
  \item $(\ip{x^k, y^k})\to \ip{x,y}$ 
\end{itemize}
\end{lemma}
\begin{proof}
    Demostramos la última propiedad, para esta asumimos que la norma está inducida por el producto interno 
    \begin{align*}
      0\leq|\ip{x^k,y^k}-\ip{x,y}|=\\
      |\ip{x^k,y^k}-\ip{x^k,y}+\ip{x^k,y}-\ip{x,y}|\leq\\
      |\ip{x^k,y^k-y}|+|\ip{x^k-x,y}|\leq\\
      \norm{x^k}\norm{y^k-y}+\norm{y}\norm{x^k-x},
    \end{align*}
    ahora tomando límites y usando el toerema del Sandwich concluimos que $\lim_{k\to\infty} |\ip{x^k,y^k}-\ip{x,y}|=0$ 
\end{proof}
\begin{lemma}
  En $\R^n$ una sucesion converge, $(x^n)\to x$, si y solo si sus componentes convergen, $(x^k_j)\to x_j$ para $1\leq j\leq n$.
\end{lemma}
\begin{proof}
  Por el \fullref{thm:equivalencia_de_normas_en_Rn} existen $C_1,C_2>0$ tal que $C_1\norm{x^k-x}_\infty\le \norm{x^k-x}\le C_2\norm{x^k-x}_\infty$, entonces para $i\in \{1,\ldots,n\} $ \[
  C_1|x_i^k-x_i| \le C_1\max_{1\leq j\leq n}\{|x_j^k-x_j|\} \le \norm{x^k-x}\le C_2\max_{1\leq j\leq n}\{|x_j^k-x_j|\}.
  \]
  Asume que $(x^k)\to x$; sea $\epsilon>0$, existe $n_0\in \N$ tal que si $k \ge n_0$, $\norm{x^k-x} \le C_1\epsilon$ y por la cadena de desigualdades anterior esto implica que $(x_i^k)\to x_i$.\\
  Asume que para todo $i\in \{1,\ldots,n\}$, $(x_i^k)\to x_i$; sea $\epsilon >0$, existen $n_0^i \in \N$ tal que $|x_i^k-x_i|<\frac{\epsilon }{C_2}$ para $1 \leq i \leq n$; para $k \geq n_0$, donde $n_0=\max_{i} \{n_0^i\}$, tenemos que $\max_{j} \{|x_j^k-x_j\} \leq \frac{\epsilon }{C_2}$ y por la cadena de desigualdades anterior esto implica que $(x^k)\to x$ 
\end{proof}
